\begin{definition}[Composite loops of the action]
\label{def:c11-Xi}
For $x\in\{\Delta,P_b\}$ and $P=P(\varepsilon,T)$, the composite loops of $I$ along $x$, written $\Xi_x(\varepsilon)$, are the amputated one-particle-irreducible cycles of vertices of $I$ and edges $P$ in which one edge is the composite kernel along $x$, all other legs open: $\tfrac\hbar2\,\partial_{\text{composite kernel along }x}W(P,I)_{\rm tree}$ is the tree-level flow of $\Xi_x(\varepsilon)$. $\Xi^\sigma_x(\varepsilon)$ is the sum of the same loops in which one edge of the cycle other than the composite kernel is the special edge, with the sign of Proposition \ref{prop:c11-master}, minus the loops with one edge $K_\varepsilon-K_0$. The renormalised composite functional is $\Xi^{\rm ren}_x:=\lim_{\varepsilon\to0}\big(\Xi_x(\varepsilon)-I^{CT}_{\mathfrak{psl}(4|4)}[\varepsilon]^{(1)}(x)\big)$, a functional of degree $-1$ at generic momenta, whose quadratic part along $\Delta$ is the renormalised composite two-point functional of Remark \ref{rem:c9-Pi}. The composite graphs of a local functional $O$ (Definition \ref{def:c11-graphs}) are of the first type when the cycle passes through two letters of $O$ and of the second type when it passes through vertices of $I$ alone; the second type is the tree insertion of $\Xi_x(\varepsilon)$ into $O$, written $\mathrm{tr}(\Xi_x,O)$.
\end{definition}

\begin{lemma}[Images of the composite loops]
\label{lem:c11-Cid}
For every $\varepsilon>0$, as functionals at generic momenta, $s\,\Xi_x(\varepsilon)=\Xi^\sigma_x(\varepsilon)$. Consequently $\Xi^\sigma_x$ has no term in $\log\varepsilon$ and no poles, $I^{CT}_{\mathfrak{psl}(4|4)}[\varepsilon]^{(1)}(x)=\log\varepsilon\,L^{(1)}(x)$ with $L^{(1)}(x)$ an $s$-closed local functional, and the composite part along $x$ is $s\,\Xi^{\rm ren}_x$, the coefficient of $\varepsilon^0$ of $\Xi^\sigma_x$, a local functional.
\end{lemma}

\begin{proof}
Compute $Q$ of the tree-level flow of $\Xi_x$ in two ways. First, $\tfrac\hbar2\,\partial_{\text{composite kernel}}W(P,I)_{\rm tree}$ is that flow, and $\partial$ of the composite kernel anticommutes with $Q$ (Lemma \ref{lem:c11-kernels} (v)); by Proposition \ref{prop:c4-defectflow} at order $\hbar^0$ and the classical master equation, $Q$ of it is $-\tfrac\hbar2\,\partial$ of the composite kernel applied to the tree-level flow of $\tfrac12\{I,I\}_\varepsilon-\tfrac12\{I,I\}$, the composite closures of the trees of $I$ with one special edge (Proposition \ref{prop:c11-master} with $M=0$). Sort them by the position of the special edge: on the cycle closed by the composite kernel, which gives the tree-level flow of $\Xi^\sigma_x$, or on a tree, between two vertices of $I$ or between $\Xi_x$ and a vertex of $I$. Second, with $\Xi_x$ as a vertex, Proposition \ref{prop:c4-defectflow} applied to $I+\delta\,\Xi_x$, $\delta$ of degree $+1$, gives the tree-level flow of $s\,\Xi_x$ plus the tree-level flows with the special-edge parts of $\{I,\Xi_x\}_\varepsilon$ and of $\tfrac12\{I,I\}_\varepsilon$ on tree edges. The graphs with the special edge on a tree are the same in both evaluations, with the same signs, the sign rule of Proposition \ref{prop:c11-master} being the Koszul rule of the second evaluation. Hence the tree-level flows of $s\,\Xi_x$ and of $\Xi^\sigma_x$ agree, and the tree-level flow is injective on functionals graded by the number of legs (the argument of Theorem \ref{thm:c8-L} (i)). $\Xi^\sigma_x$ is a sum of cycles with one special edge and the composite kernel, which by the argument of Theorem \ref{thm:c11-simplex} are $\sum_m\varepsilon^{-m}a_m+a_0+o(1)$ without $\log\varepsilon$. The poles of $I^{CT}_{\mathfrak{psl}(4|4)}[\varepsilon]^{(1)}(x)$ vanish by the weight count of Proposition \ref{prop:c9-structure} (i), and along $P_b$ by the same count at Weyl weight $5-2m$; so $I^{CT}_{\mathfrak{psl}(4|4)}[\varepsilon]^{(1)}(x)=\log\varepsilon\,L^{(1)}(x)$, and the coefficient of $\log\varepsilon$ in $s\,\Xi_x=s\,I^{CT}_{\mathfrak{psl}(4|4)}[\varepsilon]^{(1)}(x)+s\,\Xi^{\rm ren}_x+o(1)$ gives $s\,L^{(1)}(x)=0$, while the poles of $\Xi^\sigma_x$ vanish. The coefficient of $\varepsilon^0$ gives the last statement.
\end{proof}

\begin{lemma}[Special edges on trees of the composite graphs]
\label{lem:c11-trees}
Let $O$ be a local functional. The composite graphs of the first type with the special edge on a tree, at an output of the cycle or at a spectator letter of $O$, together with the special-edge parts of the counterterm trees of the arity-one counterterm of $O$, are $e^{-\varepsilon|p|^2}-1$ at a generic momentum $p$ times renormalised composite cycles, and tend to $0$. The composite graphs of the second type with the special edge on a tree tend to $0$, and so do the special-edge parts of the counterterm trees of $I^{CT}_{\mathfrak{psl}(4|4)}[\varepsilon]^{(1)}(x)$.
\end{lemma}

\begin{proof}
A special edge on a tree edge of momentum $p$ factorises: the graph is $(e^{-\varepsilon|p|^2}-1)$ times the composite cycle through the pair with its trees; the counterterm vertex joined by the same special edge gives $(e^{-\varepsilon|p|^2}-1)$ times minus its singular part, so the sum is $(e^{-\varepsilon|p|^2}-1)$ times the renormalised composite cycle, which has a limit at generic momenta (Proposition \ref{prop:c9-structure} (i)). At a spectator letter of $O$ the same holds with the special-edge parts of $\{I,O\}_\varepsilon$ and of the bracket of $I$ with the arity-one counterterm of $O$. The graphs of the second type are $(e^{-\varepsilon|p|^2}-1)$ times $\Xi_x(\varepsilon)$, which is $\log\varepsilon\,L^{(1)}(x)$ plus a convergent functional (Lemma \ref{lem:c11-Cid}), and $(e^{-\varepsilon|p|^2}-1)\log\varepsilon\to0$; the counterterm trees carry $(e^{-\varepsilon|p|^2}-1)\log\varepsilon$.
\end{proof}

\begin{proposition}[The image of the composite graphs of each type]
\label{prop:c11-typeI}
For every local functional $O$ and every $\varepsilon>0$, with (A) the special-edge composite cycles of the pairs of letters of $O$ with trees of $I$ on their outputs (Definition \ref{def:c11-sp}), (B) the composite graphs of the first type with the special edge on a tree, (D) those of the second type with the special edge on a tree, and $\epsilon_\Xi\in\{\pm1\}$ the sign with which the bracket $\{\Xi_x,O\}$ enters $Q\,\mathrm{tr}(\Xi_x,O)$,
\[
\begin{aligned}
Q\big(\text{composite graphs of the first type of }O\big)&=(\mathrm A)+(\mathrm B)-\epsilon_\Xi\{\Xi_x(\varepsilon),O\}-\big(\text{composite graphs of }sO\big)\mp\mathrm{tr}(\Xi_x,sO),\\
Q\,\mathrm{tr}(\Xi_x,O)&=\mathrm{tr}(s\,\Xi_x,O)+\epsilon_\Xi\{\Xi_x,O\}\pm\mathrm{tr}(\Xi_x,sO)+(\mathrm D),
\end{aligned}
\]
the terms with $sO$ being present only for $sO\neq0$.
\end{proposition}

\begin{proof}
By Proposition \ref{prop:c11-master} the image under $Q$ of all composite graphs of $O$ is $-\tfrac\hbar2\,\partial$ of the composite kernel applied to the trees with one vertex $sO$ or one special edge; the trees with one vertex $sO$ closed by the composite kernel are the composite graphs of $sO$, of both types. The special edge lies on the cycle through two letters (A), on a tree of a graph of the first type (B), on the cycle of a composite loop of $I$, which gives $\mathrm{tr}(\Xi^\sigma_x,O)$, or on a tree of a graph of the second type (D). The image of the second type is computed with $\Xi_x$ as a vertex, by Proposition \ref{prop:c4-defectflow} applied to $I+\delta_1\Xi_x+\delta_2O$, the part linear in $\delta_1$ and $\delta_2$: the vertex $s\,\Xi_x$ with $O$, the vertex $sO$ with $\Xi_x$, the direct contraction $\{\Xi_x,O\}$, and the special edges on trees, which are (D). Subtract: $\mathrm{tr}(s\,\Xi_x,O)=\mathrm{tr}(\Xi^\sigma_x,O)$ by Lemma \ref{lem:c11-Cid}, so the loops with the special edge on the cycle cancel, (D) cancels, and $\pm\mathrm{tr}(\Xi_x,sO)$ moves to the first line with the opposite sign.
\end{proof}
